% -*- coding: utf-8 -*-
% LuaLaTeX 用ドキュメント
% コンパイル: lualatex PolyserialCorrelation.tex

\documentclass[a4paper,11pt]{ltjsarticle}

% LuaTeX-ja 設定
\usepackage{luatexja}
\usepackage{luatexja-fontspec}

% 数学関連
\usepackage{amsmath,amssymb,amsthm}
\usepackage{mathtools}

% 図式・TikZ
\usepackage{tikz}
\usetikzlibrary{arrows.meta,positioning,calc,decorations.pathreplacing,shapes.geometric,patterns}
\usepackage{tikz-cd}

% ハイパーリンク
\usepackage{hyperref}
\hypersetup{
    colorlinks=true,
    linkcolor=blue!70!black,
    urlcolor=blue!50!black
}

% コード表示
\usepackage{listings}
\lstset{
    basicstyle=\ttfamily\small,
    keywordstyle=\color{blue!80!black}\bfseries,
    commentstyle=\color{green!50!black}\itshape,
    stringstyle=\color{red!70!black},
    frame=single,
    breaklines=true,
    numbers=left,
    numberstyle=\tiny\color{gray},
    xleftmargin=2em,
    framexleftmargin=1.5em,
    backgroundcolor=\color{gray!5}
}

% 色
\usepackage{xcolor}

% 定理環境
\theoremstyle{definition}
\newtheorem{definition}{定義}[section]
\newtheorem{example}[definition]{例}

\theoremstyle{plain}
\newtheorem{theorem}[definition]{定理}
\newtheorem{proposition}[definition]{命題}
\newtheorem{lemma}[definition]{補題}
\newtheorem{corollary}[definition]{系}

\theoremstyle{remark}
\newtheorem{remark}[definition]{注意}
\newtheorem{insight}[definition]{知見}

% カスタムコマンド
\newcommand{\R}{\mathbb{R}}
\newcommand{\N}{\mathbb{N}}
\newcommand{\Phi}{\Phi}
\newcommand{\CDF}{\mathrm{CDF}}
\newcommand{\PDF}{\mathrm{PDF}}
\newcommand{\Var}{\mathrm{Var}}
\newcommand{\E}{\mathbb{E}}

% タイトル
\title{%
    \textbf{ポリシリアル相関の数学的基礎}\\[0.5em]
    \Large Mathematical Foundations of Polyserial Correlation
}

\author{su}

\date{\today}

\begin{document}

\maketitle

\begin{center}
\small
\textit{AI assistance disclosure:}
本 \TeX 文書および Lean ソースコードは AI ツール（Antigravity (Google DeepMind), Codex (OpenAI)）によって生成され、著者による加筆・修正は行っていない。内容の正確性は保証されず、誤りがあれば著者の責任である。
\end{center}

\begin{abstract}
本文書は、連続変数と順序カテゴリ変数の間の相関を測る「ポリシリアル相関」の数学的基礎を
学部生向けに解説する。潜在変数モデルに基づき、観測されるカテゴリの背後に連続的な
潜在変数が存在すると仮定し、最尤推定による相関係数の推定手法を定式化する。
本内容は Lean 4 で形式化されている。
\end{abstract}

\tableofcontents

\newpage

%=============================================================================
\section{導入}
%=============================================================================

\subsection{ポリシリアル相関とは}

統計学において、変数間の相関を調べることは基礎的かつ重要な分析である。
連続変数同士の相関にはピアソンの相関係数が用いられるが、
一方の変数が「順序カテゴリ変数」（例：満足度を「低・中・高」で評価）の場合はどうだろうか？

\begin{definition}[順序カテゴリ変数]
\textbf{順序カテゴリ変数}（ordinal variable）とは、$k$ 個のカテゴリを持ち、
カテゴリ間に自然な順序が存在する変数である。例えば：
\begin{itemize}
    \item 満足度：「不満」$<$「普通」$<$「満足」
    \item 教育年数：「中卒」$<$「高卒」$<$「大卒」$<$「大学院卒」
    \item 評価：$1, 2, 3, 4, 5$ 点
\end{itemize}
\end{definition}

\begin{definition}[ポリシリアル相関]
\textbf{ポリシリアル相関}（polyserial correlation）は、
連続変数 $X$ と順序カテゴリ変数 $Y$（$k \geq 2$ カテゴリ）の間の
相関を推定する手法である。
\end{definition}

\begin{figure}[htbp]
\centering
\begin{tikzpicture}[
    scale=1.0,
    box/.style={draw=blue!70, fill=blue!10, rounded corners, minimum width=2.5cm, minimum height=1cm, align=center},
    arrow/.style={-Stealth, thick},
    label/.style={font=\small}
]

% 変数の種類
\node[box] (cont1) at (-4, 2) {連続変数\\$X_1$};
\node[box] (cont2) at (-4, 0) {連続変数\\$X_2$};
\node[box] (ord) at (0, 1) {順序カテゴリ\\$Y$};
\node[box] (cont3) at (4, 1) {連続変数\\$X$};

% 相関の種類
\draw[arrow, green!50!black] (cont1.south) -- (cont2.north) 
    node[midway, left, label] {ピアソン相関};

\draw[arrow, red!70] (ord.east) -- (cont3.west) 
    node[midway, above, label] {ポリシリアル相関};

% 背景
\node[label, gray] at (0, -1) {本文書で扱う};

\end{tikzpicture}
\caption{変数の種類と相関係数の選択}
\label{fig:correlation-types}
\end{figure}

\subsection{なぜ単純なピアソン相関ではダメか}

順序カテゴリ変数に単純に数値（例：$1, 2, 3$）を割り当ててピアソン相関を計算すると、
以下の問題が生じる：

\begin{itemize}
    \item カテゴリ間の「距離」が等間隔であると暗黙に仮定してしまう
    \item 相関係数が過小評価される傾向がある
    \item 統計的推論の妥当性が失われる
\end{itemize}

ポリシリアル相関は、これらの問題を「潜在変数モデル」によって解決する。

%=============================================================================
\section{標準正規分布の復習}
%=============================================================================

ポリシリアル相関を理解するためには、正規分布の基本事項を押さえておく必要がある。

\subsection{標準正規分布}

\begin{definition}[標準正規分布のPDF]
標準正規分布の\textbf{確率密度関数}（PDF）は次で定義される：
\[
    \phi(x) = \frac{1}{\sqrt{2\pi}} \exp\left(-\frac{x^2}{2}\right)
\]
\end{definition}

\begin{definition}[標準正規分布のCDF]
標準正規分布の\textbf{累積分布関数}（CDF）は次で定義される：
\[
    \Phi(x) = \int_{-\infty}^{x} \phi(t)\, dt = \int_{-\infty}^{x} \frac{1}{\sqrt{2\pi}} \exp\left(-\frac{t^2}{2}\right) dt
\]
\end{definition}

\begin{figure}[htbp]
\centering
\begin{tikzpicture}[
    scale=0.9,
    domain=-3.5:3.5,
    samples=100,
    axis/.style={-Stealth, thick}
]

% PDF のグラフ
\begin{scope}[xshift=-5cm]
    \draw[axis] (-3.5, 0) -- (3.5, 0) node[right] {$x$};
    \draw[axis] (0, -0.2) -- (0, 2) node[above] {$\phi(x)$};
    
    \draw[thick, blue!70, domain=-3.3:3.3] 
        plot (\x, {1.5*exp(-\x*\x/2)});
    
    % ピーク
    \draw[dashed, gray] (0, 0) -- (0, 1.5);
    \node[below, font=\small] at (0, -0.3) {$0$};
    \node[left, font=\small] at (-0.1, 1.5) {$\frac{1}{\sqrt{2\pi}}$};
    
    \node[font=\small, blue!70!black] at (0, -1) {確率密度関数 (PDF)};
\end{scope}

% CDF のグラフ
\begin{scope}[xshift=5cm]
    \draw[axis] (-3.5, 0) -- (3.5, 0) node[right] {$x$};
    \draw[axis] (0, -0.2) -- (0, 2) node[above] {$\Phi(x)$};
    
    % 近似的なCDFの形状
    \draw[thick, red!70, domain=-3.3:3.3] 
        plot (\x, {1.5/(1 + exp(-1.7*\x))});
    
    % 水平線
    \draw[dashed, gray] (-3.5, 1.5) -- (3.5, 1.5);
    \node[left, font=\small] at (-3.5, 1.5) {$1$};
    \node[left, font=\small] at (-3.5, 0.75) {$0.5$};
    \draw[dashed, gray] (0, 0) -- (0, 0.75);
    
    \node[font=\small, red!70!black] at (0, -1) {累積分布関数 (CDF)};
\end{scope}

\end{tikzpicture}
\caption{標準正規分布の PDF と CDF}
\label{fig:normal-dist}
\end{figure}

\begin{theorem}[正規分布の基本性質]
以下の性質が成り立つ：
\begin{enumerate}
    \item $\phi(x) \geq 0$ （非負性）
    \item $\displaystyle\int_{-\infty}^{\infty} \phi(x)\, dx = 1$ （正規化）
    \item $\Phi$ は単調増加
    \item $\displaystyle\lim_{x \to -\infty} \Phi(x) = 0$, $\displaystyle\lim_{x \to +\infty} \Phi(x) = 1$
\end{enumerate}
\end{theorem}

\subsection{二変量正規分布}

\begin{definition}[二変量標準正規分布]
二変量標準正規分布（相関係数 $\rho$, $|\rho| < 1$）の確率密度関数は：
\[
    f(x, y; \rho) = \frac{1}{2\pi\sqrt{1-\rho^2}} 
    \exp\left(-\frac{x^2 - 2\rho xy + y^2}{2(1-\rho^2)}\right)
\]
\end{definition}

\begin{proposition}[条件付き分布]
$(X, Y^*) \sim \text{BVN}(0, 0, 1, 1, \rho)$%
\footnote{$\text{BVN}(\mu_X, \mu_{Y^*}, \sigma_X^2, \sigma_{Y^*}^2, \rho)$ は二変量正規分布を表し、
ここでは平均 $(0, 0)$、分散 $(1, 1)$、相関 $\rho$ を意味する。}
のとき、$X = x$ を与えたときの $Y^*$ の条件付き分布は：
\begin{align}
    \E[Y^* \mid X = x] &= \rho x \\
    \Var[Y^* \mid X = x] &= 1 - \rho^2
\end{align}
\end{proposition}

この条件付き分布の性質が、ポリシリアル相関の根幹をなす。

%=============================================================================
\section{潜在変数モデル}
%=============================================================================

\subsection{モデルの直感}

ポリシリアル相関の核となるアイデアは、観測される順序カテゴリ変数 $Y$ の背後に
「連続的な潜在変数 $Y^*$」が存在すると仮定することである。

\begin{figure}[htbp]
\centering
\begin{tikzpicture}[
    scale=1.0,
    axis/.style={-Stealth, thick},
    threshold/.style={thick, red!70, dashed}
]

% 潜在変数の軸
\draw[axis] (-1, 0) -- (8, 0) node[right] {$Y^*$（潜在変数）};

% 閾値
\draw[threshold] (2, -0.5) -- (2, 2.5);
\draw[threshold] (5, -0.5) -- (5, 2.5);

% 閾値のラベル
\node[below, font=\small, red!70] at (2, -0.5) {$\tau_1$};
\node[below, font=\small, red!70] at (5, -0.5) {$\tau_2$};

% カテゴリ
\node[font=\small, blue!70!black, align=center] at (0.5, 1.5) {$Y = 0$\\「低」};
\node[font=\small, blue!70!black, align=center] at (3.5, 1.5) {$Y = 1$\\「中」};
\node[font=\small, blue!70!black, align=center] at (6.5, 1.5) {$Y = 2$\\「高」};

% 矢印
\draw[thick, blue!50, decorate, decoration={brace, amplitude=5pt}] 
    (-0.5, 2.3) -- (1.9, 2.3) node[midway, above=5pt, font=\small] {$Y^* < \tau_1$};
\draw[thick, blue!50, decorate, decoration={brace, amplitude=5pt}] 
    (2.1, 2.3) -- (4.9, 2.3) node[midway, above=5pt, font=\small] {$\tau_1 \leq Y^* < \tau_2$};
\draw[thick, blue!50, decorate, decoration={brace, amplitude=5pt}] 
    (5.1, 2.3) -- (7.5, 2.3) node[midway, above=5pt, font=\small] {$Y^* \geq \tau_2$};

% 境界条件
\node[font=\scriptsize, gray] at (-0.5, -1) {$\tau_0 = -\infty$};
\node[font=\scriptsize, gray] at (7.5, -1) {$\tau_3 = +\infty$};

\end{tikzpicture}
\caption{潜在変数モデル：連続的な $Y^*$ が閾値によってカテゴリ化される}
\label{fig:latent-variable}
\end{figure}

\subsection{形式的な定義}

\begin{definition}[閾値パラメータ]
$k$ カテゴリの順序変数に対して、閾値パラメータ $\boldsymbol{\tau} = (\tau_1, \ldots, \tau_{k-1})$ は
狭義単調増加：
\[
    -\infty = \tau_0 < \tau_1 < \tau_2 < \cdots < \tau_{k-1} < \tau_k = +\infty
\]
を満たす実数列である。
\end{definition}

\begin{definition}[潜在変数モデル]
潜在変数モデルは以下の構造を持つ：
\begin{itemize}
    \item 連続変数 $X$ と潜在変数 $Y^*$ は二変量正規分布に従う
    \item 相関係数 $\rho \in (-1, 1)$
    \item 観測される順序カテゴリ $Y$ は $Y^*$ を閾値で離散化したもの：
    \[
        Y = j \quad \Longleftrightarrow \quad \tau_j \leq Y^* < \tau_{j+1}
    \]
\end{itemize}
\end{definition}

\subsection{条件付きカテゴリ確率}

\begin{definition}[カテゴリ確率]
連続変数 $X = x$ を与えたとき、順序カテゴリ $Y = j$ となる確率は：
\[
    P(Y = j \mid X = x) = \Phi\left(\frac{\tau_{j+1} - \rho x}{\sigma}\right) 
                        - \Phi\left(\frac{\tau_j - \rho x}{\sigma}\right)
\]
ただし $\sigma = \sqrt{1 - \rho^2}$ である。

\textbf{注}：$j \in \{0, 1, \ldots, k-1\}$ であり、Lean 4 では \texttt{Fin k} で表現される。
\end{definition}

\begin{figure}[htbp]
\centering
\begin{tikzpicture}[
    scale=0.85,
    axis/.style={-Stealth, thick}
]

% 条件付き正規分布
\draw[axis] (-1, 0) -- (8, 0) node[right] {$Y^*$};
\draw[axis] (0, -0.2) -- (0, 2.5) node[above] {密度};

% 分布曲線
\draw[thick, blue!70, domain=-0.5:7.5, samples=100] 
    plot (\x, {2*exp(-(\x-3.5)*(\x-3.5)/3)});

% 閾値
\draw[dashed, red!70] (2, 0) -- (2, 2.2);
\draw[dashed, red!70] (5, 0) -- (5, 2.2);
\node[below, font=\small, red!70] at (2, 0) {$\tau_1$};
\node[below, font=\small, red!70] at (5, 0) {$\tau_2$};

% 塗りつぶし（カテゴリ1）
\fill[green!30, opacity=0.7, domain=2:5, samples=50] 
    (2, 0) -- plot (\x, {2*exp(-(\x-3.5)*(\x-3.5)/3)}) -- (5, 0) -- cycle;

% ラベル
\node[font=\small, green!50!black] at (3.5, 0.8) {$P(Y=1 \mid X=x)$};

% 平均の位置
\draw[dashed, gray] (3.5, 0) -- (3.5, 2);
\node[below, font=\small] at (3.5, -0.3) {$\rho x$};

\end{tikzpicture}
\caption{条件付きカテゴリ確率：条件付き分布 $Y^* \mid X = x \sim N(\rho x, 1-\rho^2)$ における閾値間の面積がカテゴリ確率を与える}
\label{fig:category-prob}
\end{figure}

\begin{theorem}[確率の正規化]
カテゴリ確率は正規化されている：
\[
    \sum_{j=0}^{k-1} P(Y = j \mid X = x) = 1
\]
\end{theorem}

\begin{proof}[証明の概略]
CDFの性質と閾値の境界条件 $\tau_0 = -\infty$, $\tau_k = +\infty$ より、
\begin{align*}
    \sum_{j=0}^{k-1} P(Y = j \mid X = x) 
    &= \sum_{j=0}^{k-1} \left[\Phi\left(\frac{\tau_{j+1} - \rho x}{\sigma}\right) 
        - \Phi\left(\frac{\tau_j - \rho x}{\sigma}\right)\right] \\
    &= \Phi(+\infty) - \Phi(-\infty) = 1 - 0 = 1
\end{align*}
（望遠鏡和の原理による）
\end{proof}

%=============================================================================
\section{最尤推定}
%=============================================================================

\subsection{尤度関数}

$n$ 個の観測データ $(x_1, y_1), \ldots, (x_n, y_n)$ が与えられたとき、
パラメータ $(\rho, \boldsymbol{\tau})$ の尤度関数を構成する。

\begin{definition}[単一観測の尤度]
観測 $(x, y)$ に対する尤度は：
\[
    L(\rho, \boldsymbol{\tau}; x, y) = P(Y = y \mid X = x) \cdot \phi(x)
\]
ここで $\phi(x)$ は $X$ の周辺密度である。
\end{definition}

\begin{definition}[対数尤度関数]
全データに対する対数尤度は：
\[
    \ell(\rho, \boldsymbol{\tau}) = \sum_{i=1}^{n} \log L(\rho, \boldsymbol{\tau}; x_i, y_i)
\]
\end{definition}

\subsection{ポリシリアル相関の推定}

\begin{definition}[ポリシリアル相関係数]
ポリシリアル相関係数 $\hat{\rho}$ は、対数尤度を最大化するパラメータとして定義される：
\[
    \hat{\rho} = \arg\max_{\rho \in (-1, 1)} \max_{\boldsymbol{\tau}} \ell(\rho, \boldsymbol{\tau})
\]
\end{definition}

\begin{figure}[htbp]
\centering
\begin{tikzpicture}[
    scale=1.0,
    box/.style={draw=blue!70, fill=blue!10, rounded corners, minimum width=3cm, minimum height=1.2cm, align=center},
    arrow/.style={-Stealth, thick}
]

% フロー
\node[box] (data) at (0, 0) {データ\\$(x_i, y_i)_{i=1}^n$};
\node[box] (likelihood) at (4, 0) {尤度関数\\$L(\rho, \boldsymbol{\tau})$};
\node[box] (optimize) at (8, 0) {最大化\\$\arg\max$};
\node[box] (result) at (12, 0) {$\hat{\rho}$\\ポリシリアル相関};

\draw[arrow] (data) -- (likelihood);
\draw[arrow] (likelihood) -- (optimize);
\draw[arrow] (optimize) -- (result);

\end{tikzpicture}
\caption{ポリシリアル相関の推定ワークフロー}
\label{fig:estimation-flow}
\end{figure}

%=============================================================================
\section{推定量の統計的性質}
%=============================================================================

\subsection{一致性}

\begin{theorem}[一致性]
サンプルサイズ $n \to \infty$ のとき、ポリシリアル相関の推定量 $\hat{\rho}_n$ は
真の相関係数 $\rho$ に確率収束する：
\[
    \hat{\rho}_n \xrightarrow{P} \rho \quad (n \to \infty)
\]
\end{theorem}

\begin{remark}[Lean形式化との対応]
上記は確率論的な表現である。Lean4形式化では現在、決定論的バージョン
$\forall \varepsilon > 0,\ \exists N,\ \forall n \geq N,\ |\hat{\rho}_n - \rho| < \varepsilon$
を採用しており、将来的に測度論に基づく確率収束の定式化への拡張を予定している。
\end{remark}

\subsection{漸近正規性}

\begin{theorem}[漸近正規性]
適切な正則条件のもとで、
\[
    \sqrt{n}(\hat{\rho}_n - \rho) \xrightarrow{d} \mathcal{N}(0, I(\rho)^{-1})
\]
ここで $I(\rho)$ はフィッシャー情報量である。
\end{theorem}

\begin{definition}[フィッシャー情報量]
\[
    I(\rho) = \E\left[\left(\frac{\partial}{\partial \rho} \log L\right)^2\right]
\]
これは推定量の分散の逆数と関連し、$I(\rho) > 0$ である。
\end{definition}

%=============================================================================
\section{特殊ケース：二系列相関}
%=============================================================================

\begin{theorem}[$k = 2$ の場合]
カテゴリ数 $k = 2$ の場合、ポリシリアル相関は\textbf{二系列相関}（biserial correlation）と一致する：
\[
    \hat{\rho}_{\text{polyserial}} = \hat{\rho}_{\text{biserial}}
\]
\end{theorem}

\begin{figure}[htbp]
\centering
\begin{tikzpicture}[
    scale=1.0,
    box/.style={draw=blue!70, fill=blue!10, rounded corners, minimum width=3.5cm, minimum height=1.5cm, align=center},
    special/.style={draw=red!70, fill=red!10, rounded corners, minimum width=3.5cm, minimum height=1.5cm, align=center}
]

% 一般ケース
\node[box] (general) at (0, 0) {ポリシリアル相関\\$k \geq 2$ カテゴリ};

% 特殊ケース
\node[special] (biserial) at (5, 0) {二系列相関\\$k = 2$ カテゴリ};

% 矢印
\draw[-Stealth, thick] (general.east) -- (biserial.west) 
    node[midway, above, font=\small] {$k = 2$};

% 例
\node[font=\small, gray, align=center] at (0, -1.5) {5段階評価\\4カテゴリ以上};
\node[font=\small, gray, align=center] at (5, -1.5) {合格/不合格\\はい/いいえ};

\end{tikzpicture}
\caption{ポリシリアル相関と二系列相関の関係}
\label{fig:biserial}
\end{figure}

%=============================================================================
\section{具体例}
%=============================================================================

\subsection{3カテゴリの例}

\begin{example}[満足度調査]
顧客満足度を「低」「中」「高」の3カテゴリで測定し、
購入金額（連続変数）との相関を調べる。

\textbf{設定}：
\begin{itemize}
    \item $k = 3$（カテゴリ数）
    \item 閾値：$\tau_1 = -1$, $\tau_2 = 1$
    \item 相関係数：$\rho = 0.5$
\end{itemize}

この設定では：
\begin{itemize}
    \item $Y = 0$（低）$\Leftrightarrow$ $Y^* < -1$
    \item $Y = 1$（中）$\Leftrightarrow$ $-1 \leq Y^* < 1$
    \item $Y = 2$（高）$\Leftrightarrow$ $Y^* \geq 1$
\end{itemize}
\end{example}

\begin{figure}[htbp]
\centering
\begin{tikzpicture}[
    scale=0.9,
    axis/.style={-Stealth, thick}
]

% 座標軸
\draw[axis] (-0.5, 0) -- (6, 0) node[right] {$X$（購入金額）};
\draw[axis] (0, -0.5) -- (0, 4) node[above] {$Y$（満足度）};

% データ点（模擬）
\foreach \x/\y in {0.5/0, 1/0, 1.2/0, 0.8/0, 1.5/1, 2/1, 2.2/1, 2.5/1, 2.8/1, 3/1, 
                   3.2/2, 3.5/2, 4/2, 4.5/2, 5/2, 5.2/2} {
    \filldraw[blue!70, opacity=0.7] (\x, \y+0.5) circle (3pt);
}

% カテゴリラベル
\node[left, font=\small] at (0, 0.5) {低 (0)};
\node[left, font=\small] at (0, 1.5) {中 (1)};
\node[left, font=\small] at (0, 2.5) {高 (2)};

% 水平線
\draw[dashed, gray] (0, 1) -- (6, 1);
\draw[dashed, gray] (0, 2) -- (6, 2);

% トレンド
\draw[thick, red!70, -Stealth] (0.5, 0.3) -- (5.5, 2.7) 
    node[right, font=\small] {$\rho = 0.5$};

\end{tikzpicture}
\caption{3カテゴリの順序変数と連続変数の散布図（模擬データ）}
\label{fig:example-data}
\end{figure}

%=============================================================================
\section{Lean形式化の概要}
%=============================================================================

本文書の内容は Lean 4 で形式化されている。主要な構成要素を紹介する。

\subsection{データ構造}

\begin{lstlisting}[language=ML, caption={順序カテゴリ変数の定義}]
structure OrdinalVariable (k : Nat) where
  value : Fin k
  deriving DecidableEq
\end{lstlisting}

\begin{lstlisting}[language=ML, caption={閾値パラメータの定義}]
structure Thresholds (k : Nat) where
  tau : Fin (k - 1) -> Real
  strictly_increasing : StrictMono tau
\end{lstlisting}

\begin{lstlisting}[language=ML, caption={潜在変数モデルの定義}]
structure LatentVariableModel (k : Nat) where
  thresholds : Thresholds k
  rho : Real
  h_rho : |rho| < 1
  h_k : 1 < k
\end{lstlisting}

\subsection{主要定理}

\begin{lstlisting}[language=ML, caption={確率の正規化}]
theorem categoryProbability_sum_eq_one {k : Nat} (hk : 0 < k)
    (model : LatentVariableModel k) (x : Real) :
    Finset.sum Finset.univ (fun j => categoryProbability model x j) = 1
\end{lstlisting}

\begin{lstlisting}[language=ML, caption={二系列相関との一致}]
theorem polyserial_eq_biserial_when_k_eq_2 {n : Nat}
    (data : Fin n -> Real * Fin 2) :
    polyserialCorrelation data = biserialCorrelation data
\end{lstlisting}

\subsection{形式化のアプローチ}

本形式化では、以下のように公理的アプローチを採用している：

\begin{enumerate}
    \item 標準正規分布の性質（非負性、正規化、単調性など）を公理として導入
    \item これらの公理から定理を導出
    \item 将来的に mathlib の測度論・確率論の結果で公理を置換予定
\end{enumerate}

主要な公理（\texttt{PolyserialAxioms} 名前空間）：
\begin{itemize}
    \item \texttt{Axiom\_stdNormalPDF\_nonneg} --- $\phi(x) \geq 0$
    \item \texttt{Axiom\_stdNormalPDF\_integral\_eq\_one} --- $\int \phi = 1$
    \item \texttt{Axiom\_stdNormalCDF\_mono} --- $\Phi$ の単調性
    \item \texttt{Axiom\_categoryProbability\_sum\_eq\_one} --- カテゴリ確率の正規化
    \item \texttt{Axiom\_polyserial\_consistent} --- 推定量の一致性
\end{itemize}

%=============================================================================
\section{用語対応表}
%=============================================================================

\begin{table}[htbp]
\centering
\caption{Lean識別子と数学記号の対応}
\small
\begin{tabular}{|l|l|l|}
\hline
\textbf{Lean 識別子} & \textbf{数学記号} & \textbf{説明} \\
\hline
\texttt{stdNormalPDF} & $\phi(x)$ & 標準正規分布のPDF \\
\texttt{stdNormalCDF} & $\Phi(x)$ & 標準正規分布のCDF \\
\texttt{Thresholds.tau} & $\tau_j$ & 閾値パラメータ \\
\texttt{LatentVariableModel.rho} & $\rho$ & 相関係数 \\
\texttt{categoryProbability} & $P(Y = j \mid X = x)$ & カテゴリ確率 \\
\texttt{polyserialCorrelation} & $\hat{\rho}$ & ポリシリアル相関推定量 \\
\texttt{fisherInformation} & $I(\rho)$ & フィッシャー情報量 \\
\hline
\end{tabular}
\end{table}

%=============================================================================
\section{まとめと発展}
%=============================================================================

\subsection{本文書のまとめ}

\begin{enumerate}
    \item ポリシリアル相関は、連続変数と順序カテゴリ変数の間の相関を推定する手法
    \item 潜在変数モデルに基づき、観測カテゴリの背後に連続変数を仮定
    \item 最尤推定により相関係数を推定
    \item 推定量は一致性と漸近正規性を持つ
    \item $k = 2$ の場合は二系列相関と一致
\end{enumerate}

\subsection{発展的話題}

\begin{itemize}
    \item \textbf{ポリコリック相関}：両変数が順序カテゴリの場合への一般化
    \item \textbf{欠測データ処理}：不完全なデータへの対応
    \item \textbf{ベイズ推定}：事前分布を用いた推定手法
    \item \textbf{多変量への拡張}：複数の変数間の相関構造
\end{itemize}

%=============================================================================
% 参考文献
%=============================================================================

\begin{thebibliography}{9}

\bibitem{olsson1979}
U. Olsson, F. Drasgow, and N. J. Dorans.
\textit{The polyserial correlation coefficient}.
Psychometrika, 47(3):337--347, 1982.

\bibitem{mathlib}
The mathlib Community.
\textit{Mathlib: The Math Library for Lean 4}.
\url{https://github.com/leanprover-community/mathlib4}

\bibitem{lean4}
Leonardo de Moura et al.
\textit{The Lean 4 Theorem Prover and Programming Language}.
CADE-28, 2021.

\bibitem{degenne2023}
Rémy Degenne et al.
\textit{Probability Theory in Mathlib}.
ITP 2023.
\url{https://doi.org/10.4230/LIPIcs.ITP.2023.12}

\end{thebibliography}

\end{document}
